% sec3_8.tex — Section 3.8: Implications of Conditional Homoskedasticity

So far, we have not assumed conditional homoskedasticity in developing the asymptotics
of the GMM estimator.  This section considers the implication of imposing

\begin{assumption}[conditional homoskedasticity]
\label{ass:3.7}
\[
  \E(\varepsilon_i^2 \mid \mathbf{x}_i) = \sigma^2.
\]
\end{assumption}

Under conditional homoskedasticity, the matrix of fourth moments $\mathbf{S}$
$(= \E(\mathbf{g}_i \mathbf{g}_i') = \E(\varepsilon_i^2\, \mathbf{x}_i \mathbf{x}_i'))$ can be written as a
product of second moments:
\begin{equation}
  \mathbf{S} = \sigma^2 \boldsymbol{\Sigma}_{\mathbf{xx}},
  \label{eq:3.8.1}
\end{equation}
where $\boldsymbol{\Sigma}_{\mathbf{xx}} = \E(\mathbf{x}_i \mathbf{x}_i')$.  As in Chapter~2, this
decomposition of $\mathbf{S}$ has several implications.

\begin{itemize}
\item Since $\mathbf{S}$ is nonsingular by Assumption~3.5, this decomposition implies that
  $\sigma^2 > 0$ and $\boldsymbol{\Sigma}_{\mathbf{xx}}$ is nonsingular.

\item The estimator exploiting this structure of $\mathbf{S}$ is
  \begin{equation}
    \hat{\mathbf{S}} = \hat{\sigma}^2 \cdot \frac{1}{n} \sum_{i=1}^{n}
    \mathbf{x}_i \mathbf{x}_i' = \hat{\sigma}^2\, \mathbf{S}_{\mathbf{xx}},
    \label{eq:3.8.2}
  \end{equation}
  where $\hat{\sigma}^2$ is some consistent estimator to be specified below.  By ergodic stationarity,
  $\mathbf{S}_{\mathbf{xx}} \to_{\text{a.s.}} \boldsymbol{\Sigma}_{\mathbf{xx}}$.  Thus, provided that
  $\hat{\sigma}^2$ is consistent, we do not need the fourth-moment assumption (Assumption~3.6)
  for $\hat{\mathbf{S}}$ to be consistent.
\end{itemize}

Needless to say, all the results presented so far are valid under the extra
condition of conditional homoskedasticity.  But many of the results and formulas
can be simplified substantially under this extra condition, by just replacing $\mathbf{S}$
by $\sigma^2 \boldsymbol{\Sigma}_{\mathbf{xx}}$ and the expression \eqref{eq:3.5.10} for
$\hat{\mathbf{S}}$ by \eqref{eq:3.8.2}.  This section collects those simplifications.

%% ─── Efficient GMM Becomes 2SLS ────────────────────────────────────────────
\subsection*{Efficient GMM Becomes 2SLS}

In the efficient GMM estimation, the weighting matrix is $\hat{\mathbf{S}}^{-1}$.  If we set
$\hat{\mathbf{S}}$ to \eqref{eq:3.8.2}, the GMM estimator becomes
\begin{align}
  \hat{\boldsymbol{\delta}}(\hat{\mathbf{S}}^{-1})
  &= [\mathbf{S}_{\mathbf{xz}}'(\hat{\sigma}^2\, \mathbf{S}_{\mathbf{xx}})^{-1}
     \mathbf{S}_{\mathbf{xz}}]^{-1}\,
     \mathbf{S}_{\mathbf{xz}}'(\hat{\sigma}^2\, \mathbf{S}_{\mathbf{xx}})^{-1}\,
     \mathbf{s}_{\mathbf{xy}} \notag \\
  &= (\mathbf{S}_{\mathbf{xz}}' \mathbf{S}_{\mathbf{xx}}^{-1}
     \mathbf{S}_{\mathbf{xz}})^{-1}\,
     \mathbf{S}_{\mathbf{xz}}' \mathbf{S}_{\mathbf{xx}}^{-1}\,
     \mathbf{s}_{\mathbf{xy}} \notag \\
  &= \hat{\boldsymbol{\delta}}(\mathbf{S}_{\mathbf{xx}}^{-1})
  \equiv \hat{\boldsymbol{\delta}}_{\text{2SLS}},
  \label{eq:3.8.3}
\end{align}
which does not depend on $\hat{\sigma}^2$.  In the general case, the whole point of the first step in
the efficient two-step GMM was to obtain a consistent estimator of $\mathbf{S}$.  Under conditional
homoskedasticity, there is no need to do the first step because the second-step
estimator collapses to the GMM estimator with $\mathbf{S}_{\mathbf{xx}}$ used for $\hat{\mathbf{S}}$,
$\hat{\boldsymbol{\delta}}(\mathbf{S}_{\mathbf{xx}}^{-1})$.  This
estimator, $\hat{\boldsymbol{\delta}}_{\text{2SLS}}$, is called the \textbf{Two-Stage Least Squares}
(2SLS or TSLS).\footnote{This estimator was first proposed by Theil (1953).}
The same equation can be estimated by maximum likelihood.  Section~8.6 will cover
the ML counterpart of 2SLS, called the ``limited-information maximum likelihood
estimator.''

The expression for $\avar(\hat{\boldsymbol{\delta}}_{\text{2SLS}})$ can be obtained by substituting
\eqref{eq:3.8.1} into \eqref{eq:3.5.13}:
\begin{equation}
  \avar(\hat{\boldsymbol{\delta}}_{\text{2SLS}})
  = \sigma^2 \cdot (\boldsymbol{\Sigma}_{\mathbf{xz}}'
    \boldsymbol{\Sigma}_{\mathbf{xx}}^{-1}
    \boldsymbol{\Sigma}_{\mathbf{xz}})^{-1}.
  \label{eq:3.8.4}
\end{equation}
A natural estimator of this is
\begin{equation}
  \widehat{\avar(\hat{\boldsymbol{\delta}}_{\text{2SLS}})}
  = \hat{\sigma}^2 \cdot (\mathbf{S}_{\mathbf{xz}}'
    \mathbf{S}_{\mathbf{xx}}^{-1}
    \mathbf{S}_{\mathbf{xz}})^{-1}.
  \label{eq:3.8.5}
\end{equation}
For $\hat{\sigma}^2$, consider the sample variance of the 2SLS residuals:
\begin{equation}
  \hat{\sigma}^2 \equiv \frac{1}{n} \sum_{i=1}^{n}
  (y_i - \mathbf{z}_i' \hat{\boldsymbol{\delta}}_{\text{2SLS}})^2.
  \label{eq:3.8.6}
\end{equation}
(Some authors divide the sum of squares by $n - L$, not by $n$, to calculate $\hat{\sigma}^2$.)
By Proposition~3.2, \eqref{eq:3.8.6} $\pto \sigma^2$ if
$\E(\mathbf{z}_i \mathbf{z}_i')$ exists and is finite.\footnote{This additional assumption is needed
for the consistency of $\hat{\sigma}^2$; see Proposition~3.2.}  Thus $\hat{\mathbf{S}}$ defined in
\eqref{eq:3.8.2} with this $\hat{\sigma}^2$ is consistent for $\mathbf{S}$.

Substituting \eqref{eq:3.8.2} into \eqref{eq:3.5.15} and \eqref{eq:3.5.16}, the $t$-ratio
and the Wald statistic become
\begin{align}
  t_\ell &= \frac{\hat{\delta}_{\text{2SLS},\ell} - \bar{\delta}_\ell}{SE_\ell}
  \quad \text{with} \quad
  SE_\ell = \sqrt{\frac{\hat{\sigma}^2}{n} \cdot
    \bigl((\mathbf{S}_{\mathbf{xz}}' \mathbf{S}_{\mathbf{xx}}^{-1}
    \mathbf{S}_{\mathbf{xz}})^{-1}\bigr)_{\!\ell\ell}},
  \label{eq:3.8.7} \\
  W &= n \cdot \frac{\mathbf{a}(\hat{\boldsymbol{\delta}}_{\text{2SLS}})'\,
    \bigl[\mathbf{A}(\hat{\boldsymbol{\delta}}_{\text{2SLS}})
    (\mathbf{S}_{\mathbf{xz}}' \mathbf{S}_{\mathbf{xx}}^{-1}
    \mathbf{S}_{\mathbf{xz}})^{-1}
    \mathbf{A}(\hat{\boldsymbol{\delta}}_{\text{2SLS}})'\bigr]^{-1}\,
    \mathbf{a}(\hat{\boldsymbol{\delta}}_{\text{2SLS}})}
    {\hat{\sigma}^2}.
  \label{eq:3.8.8}
\end{align}

%% ─── J Becomes Sargan's Statistic ──────────────────────────────────────────
\subsection*{$J$ Becomes Sargan's Statistic}

When $\hat{\mathbf{W}}$ is set to $(\hat{\sigma}^2 \mathbf{S}_{\mathbf{xx}})^{-1}$, the distance
defined in \eqref{eq:3.4.6} becomes
\begin{equation}
  J(\tilde{\boldsymbol{\delta}},\, (\hat{\sigma}^2 \cdot \mathbf{S}_{\mathbf{xx}})^{-1})
  = n \cdot \frac{(\mathbf{s}_{\mathbf{xy}}
    - \mathbf{S}_{\mathbf{xz}} \tilde{\boldsymbol{\delta}})'\,
    \mathbf{S}_{\mathbf{xx}}^{-1}\,
    (\mathbf{s}_{\mathbf{xy}}
    - \mathbf{S}_{\mathbf{xz}} \tilde{\boldsymbol{\delta}})}
    {\hat{\sigma}^2}.
  \label{eq:3.8.9}
\end{equation}
Proposition~3.6 then implies that the distance evaluated at the efficient GMM estimator
under conditional homoskedasticity, $\hat{\boldsymbol{\delta}}_{\text{2SLS}}$, is asymptotically chi-squared.
This distance is called \textbf{Sargan's statistic} (Sargan, 1958):
\begin{equation}
  \text{Sargan's statistic}
  = n \cdot \frac{(\mathbf{s}_{\mathbf{xy}}
    - \mathbf{S}_{\mathbf{xz}} \hat{\boldsymbol{\delta}}_{\text{2SLS}})'\,
    \mathbf{S}_{\mathbf{xx}}^{-1}\,
    (\mathbf{s}_{\mathbf{xy}}
    - \mathbf{S}_{\mathbf{xz}} \hat{\boldsymbol{\delta}}_{\text{2SLS}})}
    {\hat{\sigma}^2}.
  \label{eq:3.8.10}
\end{equation}

We summarize our results so far as

\begin{proposition}[asymptotic properties of 2SLS]
\label{prop:3.9}
\begin{enumerate}[label=(\alph*)]
\item \emph{Under Assumptions~3.1--3.4, the 2SLS estimator \eqref{eq:3.8.3} is consistent.
  If Assumption~3.5 is added, the estimator is asymptotically normal with the asymptotic
  variance given by \eqref{eq:3.5.1} with
  $\mathbf{W} = (\sigma^2 \boldsymbol{\Sigma}_{\mathbf{xx}})^{-1}$.
  If Assumption~3.7 (conditional homoskedasticity) is added to Assumptions~3.1--3.5, then
  the estimator is the efficient GMM estimator.}

\emph{Furthermore, if $\E(\mathbf{z}_i \mathbf{z}_i')$ exists and is finite, then}

\item \emph{the asymptotic variance is consistently estimated by \eqref{eq:3.8.5},}

\item $t_\ell$ \emph{in \eqref{eq:3.8.7}} $\dto N(0,1)$, $W$ \emph{in \eqref{eq:3.8.8}}
  $\dto \chi^2(\#\mathbf{r})$, \emph{and}

\item \emph{the Sargan statistic in \eqref{eq:3.8.10}} $\dto \chi^2(K - L)$.
\end{enumerate}
\end{proposition}

Proposition~3.8 states that the $LR$ statistic, which is the difference in $J$ with and
without the imposition of the null hypothesis, is asymptotically chi-squared.  Since
$J$ can be written as \eqref{eq:3.8.9}, we have
\begin{equation}
  LR = n \cdot \frac{(\mathbf{s}_{\mathbf{xy}}
    - \mathbf{S}_{\mathbf{xz}} \tilde{\boldsymbol{\delta}})'\,
    \mathbf{S}_{\mathbf{xx}}^{-1}\,
    (\mathbf{s}_{\mathbf{xy}}
    - \mathbf{S}_{\mathbf{xz}} \tilde{\boldsymbol{\delta}})
  - (\mathbf{s}_{\mathbf{xy}}
    - \mathbf{S}_{\mathbf{xz}} \hat{\boldsymbol{\delta}}_{\text{2SLS}})'\,
    \mathbf{S}_{\mathbf{xx}}^{-1}\,
    (\mathbf{s}_{\mathbf{xy}}
    - \mathbf{S}_{\mathbf{xz}} \hat{\boldsymbol{\delta}}_{\text{2SLS}})}
    {\hat{\sigma}^2},
  \label{eq:3.8.11}
\end{equation}
where $\tilde{\boldsymbol{\delta}}$ is the restricted 2SLS estimator which minimizes \eqref{eq:3.8.9}
under the null hypothesis.\footnote{This statistic was first derived by Gallant and
Jorgenson (1979).}  In Proposition~3.8, the use of the same $\hat{\mathbf{S}}$ guaranteed the
statistic to be nonnegative in finite samples.  Here, deflation by the same $\hat{\sigma}^2$
ensures the statistic to be nonnegative.  If the hypothesis is linear, then this $LR$ is
numerically equal to the Wald statistic $W$.

%% ─── Small-Sample Properties of 2SLS ──────────────────────────────────────
\subsection*{Small-Sample Properties of 2SLS}

There is a fairly large literature on the finite-sample distribution of the 2SLS estimator
(see, e.g., Judge et al.\ (1985, Section~15.4) and Staiger and Stock (1997,
Section~1)).  Some studies manage to derive analytically the exact finite-sample
distribution of the estimator, while others are Monte Carlo studies for various
DGPs.  The analytical results, however, are not useful for empirical researchers,
because they are derived under the restrictive assumptions of fixed instruments and
normal errors and the analytical expressions for distributions are computationally
intractable.

For the case of a single regressor and a single (stochastic) instrument with
normal errors, Nelson and Startz (1990) derive the exact finite-sample distribution
of the 2SLS estimator that is fairly simple and easy to calculate.  They also show
that, when the instrument is ``weak'' in the sense of low explanatory power in the
first-stage regression of the regressor on the instrument, a mass of the distribution
of the sampling error $\hat{\boldsymbol{\delta}}_{\text{2SLS}} - \boldsymbol{\delta}$ remains apart from zero
until the sample size gets really large.

Their work illustrates the need for reporting the $R^2$ for the first-stage regressions;
if the $R^2$ is low, we should suspect the large-sample approximation to the
distribution of the 2SLS estimator to be poor (you will see a dramatic example in
part~(g) of the empirical exercise).  Recently, Staiger and Stock (1997) proposed an
alternative asymptotic approximation to the finite-sample distribution of the 2SLS
and other estimators and associated test statistics, for the case of ``weak'' instruments.
Their device is to look at a sequence of models along which the coefficients
of the instruments in the first-stage regressions converge to zero.  (Analytical Exercise~10
works out this type of asymptotics for the simple case of one regressor and
one instrument.)

%% ─── Alternative Derivations of 2SLS ──────────────────────────────────────
\subsection*{Alternative Derivations of 2SLS}

If we define data matrices as
\[
  \underset{(n \times K)}{\mathbf{X}}
  = \begin{bmatrix} \mathbf{x}_1' \\ \mathbf{x}_2' \\ \vdots \\ \mathbf{x}_n' \end{bmatrix},
  \quad
  \underset{(n \times L)}{\mathbf{Z}}
  = \begin{bmatrix} \mathbf{z}_1' \\ \mathbf{z}_2' \\ \vdots \\ \mathbf{z}_n' \end{bmatrix},
  \quad
  \underset{(n \times 1)}{\mathbf{y}}
  = \begin{bmatrix} y_1 \\ y_2 \\ \vdots \\ y_n \end{bmatrix},
\]
then it is easy to see that the 2SLS estimator and associated statistics can be written as
\begin{align}
  \hat{\boldsymbol{\delta}}_{\text{2SLS}}
  &= [\mathbf{Z}'\mathbf{X}(\mathbf{X}'\mathbf{X})^{-1}\mathbf{X}'\mathbf{Z}]^{-1}\,
     \mathbf{Z}'\mathbf{X}(\mathbf{X}'\mathbf{X})^{-1}\mathbf{X}'\mathbf{y} \notag \\
  &= (\mathbf{Z}'\mathbf{P}\mathbf{Z})^{-1}\, \mathbf{Z}'\mathbf{P}\mathbf{y},
  \label{eq:3.8.3p}
\end{align}
where $\mathbf{P} \equiv \mathbf{X}(\mathbf{X}'\mathbf{X})^{-1}\mathbf{X}'$ is the projection matrix,
\begin{align}
  \widehat{\avar(\hat{\boldsymbol{\delta}}_{\text{2SLS}})}
  &= n \cdot \hat{\sigma}^2 \cdot
     [\mathbf{Z}'\mathbf{X}(\mathbf{X}'\mathbf{X})^{-1}\mathbf{X}'\mathbf{Z}]^{-1}
  = n \cdot \hat{\sigma}^2 \cdot (\mathbf{Z}'\mathbf{P}\mathbf{Z})^{-1},
  \label{eq:3.8.5p} \\
  \hat{\sigma}^2 &\equiv \frac{\hat{\boldsymbol{\varepsilon}}'\hat{\boldsymbol{\varepsilon}}}{n}
  \quad \text{where } \hat{\boldsymbol{\varepsilon}} \equiv \mathbf{y}
  - \mathbf{Z}\hat{\boldsymbol{\delta}}_{\text{2SLS}},
  \label{eq:3.8.6p} \\
  t_\ell &= \frac{\hat{\delta}_{\text{2SLS},\ell} - \bar{\delta}_\ell}
    {\sqrt{\hat{\sigma}^2 \cdot
    \bigl([\mathbf{Z}'\mathbf{X}(\mathbf{X}'\mathbf{X})^{-1}\mathbf{X}'\mathbf{Z}]^{-1}
    \bigr)_{\!\ell\ell}}},
  \label{eq:3.8.7p} \\
  W &= \frac{\mathbf{a}(\hat{\boldsymbol{\delta}}_{\text{2SLS}})'\,
    \bigl[\mathbf{A}(\hat{\boldsymbol{\delta}}_{\text{2SLS}})
    (\mathbf{Z}'\mathbf{X}(\mathbf{X}'\mathbf{X})^{-1}\mathbf{X}'\mathbf{Z})^{-1}
    \mathbf{A}(\hat{\boldsymbol{\delta}}_{\text{2SLS}})'\bigr]^{-1}\,
    \mathbf{a}(\hat{\boldsymbol{\delta}}_{\text{2SLS}})}{\hat{\sigma}^2},
  \label{eq:3.8.8p} \\
  J(\tilde{\boldsymbol{\delta}},\, (\hat{\sigma}^2 \cdot \mathbf{S}_{\mathbf{xx}})^{-1})
  &= \frac{(\mathbf{y} - \mathbf{Z}\tilde{\boldsymbol{\delta}})'\,
     \mathbf{P}\,(\mathbf{y} - \mathbf{Z}\tilde{\boldsymbol{\delta}})}{\hat{\sigma}^2},
  \label{eq:3.8.9p} \\
  \text{Sargan's statistic}
  &= \frac{\hat{\boldsymbol{\varepsilon}}'\mathbf{P}\hat{\boldsymbol{\varepsilon}}}
     {\hat{\sigma}^2}.
  \label{eq:3.8.10p}
\end{align}

Using this formula, we can provide two other derivations of the 2SLS estimator.

\medskip\noindent\textbf{2SLS as an IV Estimator.}\enspace
Let $\hat{\mathbf{z}}_i$ ($L \times 1$) be the vector of $L$ instruments (which will be generated from
$\mathbf{x}_i$ as described below) for the $L$ regressors, and let $\hat{\mathbf{Z}}$ be the $n \times L$
data matrix of those $L$ instruments.  So the $i$-th row of $\hat{\mathbf{Z}}$ is
$\hat{\mathbf{z}}_i'$.  The IV estimator of $\boldsymbol{\delta}$ with $\hat{\mathbf{z}}_i$ serving as
instruments is, by \eqref{eq:3.4.4},
\begin{equation}
  \hat{\boldsymbol{\delta}}_{\text{IV}}
  = \Bigl(\frac{1}{n} \sum_{i=1}^{n}
    \hat{\mathbf{z}}_i \mathbf{z}_i'\Bigr)^{-1}
    \frac{1}{n} \sum_{i=1}^{n} \hat{\mathbf{z}}_i \cdot y_i
  = (\hat{\mathbf{Z}}'\mathbf{Z})^{-1}\, \hat{\mathbf{Z}}'\mathbf{y}.
  \label{eq:3.8.12}
\end{equation}
Now we generate those $L$ instruments from $\mathbf{x}_i$ as follows.  The $\ell$-th instrument is
the fitted value from regressing $z_{i\ell}$ (the $\ell$-th regressor) on $\mathbf{x}_i$.  The $n$-vector
of fitted values is $\mathbf{X}(\mathbf{X}'\mathbf{X})^{-1}\mathbf{X}'\mathbf{z}_\ell$, where
$\mathbf{z}_\ell$ is the $n$-vector of the $\ell$-th regressor (i.e., the $\ell$-th column of $\mathbf{Z}$).
Therefore, the $n \times L$ data matrix of instruments is
\begin{equation}
  \hat{\mathbf{Z}}
  = (\mathbf{X}(\mathbf{X}'\mathbf{X})^{-1}\mathbf{X}'\mathbf{z}_1, \ldots,
     \mathbf{X}(\mathbf{X}'\mathbf{X})^{-1}\mathbf{X}'\mathbf{z}_L)
  = \mathbf{X}(\mathbf{X}'\mathbf{X})^{-1}\mathbf{X}'\mathbf{Z}
  = \mathbf{P}\mathbf{Z},
  \label{eq:3.8.13}
\end{equation}
where $\mathbf{P}$ is the projection matrix.  Substituting this into the IV format
\eqref{eq:3.8.12} yields the 2SLS estimator (see \eqref{eq:3.8.3p}).

\medskip\noindent\textbf{2SLS as Two Regressions.}\enspace
Instead of substituting the generated instruments $\hat{\mathbf{z}}_i$ into the IV format to estimate
the equation $y_i = \mathbf{z}_i' \boldsymbol{\delta} + \varepsilon_i$, consider regressing $y_i$ on
$\hat{\mathbf{z}}_i$.  The coefficient estimate is
\begin{align*}
  (\hat{\mathbf{Z}}'\hat{\mathbf{Z}})^{-1}\, \hat{\mathbf{Z}}'\mathbf{y}
  &= (\mathbf{Z}'\mathbf{P}'\mathbf{P}\mathbf{Z})^{-1}\,
     \mathbf{Z}'\mathbf{P}'\mathbf{P}\mathbf{y} \\
  &= (\mathbf{Z}'\mathbf{P}\mathbf{Z})^{-1}\, \mathbf{Z}'\mathbf{P}\mathbf{y}
  \quad \text{(since $\mathbf{P}$ is symmetric and idempotent)}.
\end{align*}
This is the 2SLS estimator.  So the 2SLS coefficient estimate can be obtained in two
stages: the first stage is to regress the $L$ regressors on $\mathbf{x}_i$ and obtain fitted values
$\hat{\mathbf{z}}_i$, and the second stage is to regress $y_i$ on those fitted values.

For those regressors that are predetermined and hence included in $\mathbf{x}_i$, there is no
need to carry out the first-stage regression because the regressor is the regressor
itself.  To see this, if $z_{i\ell}$ is predetermined and included in $\mathbf{x}_i$ as the $k$-th instrument,
the $n$-vector of fitted values for the $\ell$-th regressor is $\mathbf{P}\mathbf{z}_\ell$, where $\mathbf{z}_\ell$
is the $n$-vector whose $i$-th element is $z_{i\ell}$.  But since $\mathbf{z}_\ell$ is also the $k$-th column of
$\mathbf{X}$, since $\mathbf{P}$ is the projection matrix, $\mathbf{P}\mathbf{z}_\ell = \mathbf{x}_k$.
Since $\mathbf{P}$ is the projection matrix, $\mathbf{P}\mathbf{x}_k = \mathbf{x}_k$.

This derivation of the 2SLS is useful as it justifies the naming of the estimator,
but there is a pitfall.  In the second-stage regression where $y_i$ is regressed on
$\hat{\mathbf{z}}_i$, the standard errors routinely calculated by the OLS package are based on the
residual vector $\mathbf{y} - \hat{\mathbf{Z}}\hat{\boldsymbol{\delta}}_{\text{2SLS}}$.
This does \emph{not} equal the 2SLS residual
$\mathbf{y} - \mathbf{Z}\hat{\boldsymbol{\delta}}_{\text{2SLS}}$.  Therefore, the OLS standard errors and
estimated asymptotic variance from the second stage cannot be used for statistical
inference.

%% ─── When Regressors Are Predetermined ─────────────────────────────────────
\subsection*{When Regressors Are Predetermined}

When all the regressors are predetermined and the errors are conditionally homoskedastic,
there is a close connection between the distance function $J$ for efficient
GMM and the sum of squared residuals ($SSR$).  From \eqref{eq:3.8.9p} on page~230,
\begin{align}
  J(\tilde{\boldsymbol{\delta}},\, (\hat{\sigma}^2 \cdot \mathbf{S}_{\mathbf{xx}})^{-1})
  &= \frac{(\mathbf{y} - \mathbf{Z}\tilde{\boldsymbol{\delta}})'\,
     \mathbf{P}\,(\mathbf{y} - \mathbf{Z}\tilde{\boldsymbol{\delta}})}{\hat{\sigma}^2}
  \notag \\
  &= \frac{\mathbf{y}'\mathbf{P}\mathbf{y}
     - 2\mathbf{y}'\mathbf{P}\mathbf{Z}\tilde{\boldsymbol{\delta}}
     + \tilde{\boldsymbol{\delta}}'\mathbf{Z}'\mathbf{P}\mathbf{Z}\tilde{\boldsymbol{\delta}}}
     {\hat{\sigma}^2}
  \notag \\
  &= \frac{\mathbf{y}'\mathbf{P}\mathbf{y}
     - 2\mathbf{y}'\mathbf{Z}\tilde{\boldsymbol{\delta}}
     + \tilde{\boldsymbol{\delta}}'\mathbf{Z}'\mathbf{Z}\tilde{\boldsymbol{\delta}}}
     {\hat{\sigma}^2}
  \quad \text{(since $\mathbf{P}\mathbf{Z} = \mathbf{Z}$ when
    $\mathbf{z}_i \subset \mathbf{x}_i$)}
  \notag \\
  &= \frac{(\mathbf{y} - \mathbf{Z}\tilde{\boldsymbol{\delta}})'
     (\mathbf{y} - \mathbf{Z}\tilde{\boldsymbol{\delta}})}{\hat{\sigma}^2}
  - \frac{(\mathbf{y} - \hat{\mathbf{y}})'(\mathbf{y} - \hat{\mathbf{y}})}
     {\hat{\sigma}^2},
  \label{eq:3.8.14}
\end{align}
where $\hat{\mathbf{y}} \equiv \mathbf{P}\mathbf{y}$ is the vector of fitted values from unrestricted OLS.
Since the last term does not depend on $\tilde{\boldsymbol{\delta}}$, minimizing $J$ amounts to minimizing the
sum of squared residuals $(\mathbf{y} - \mathbf{Z}\tilde{\boldsymbol{\delta}})'(\mathbf{y} -
\mathbf{Z}\tilde{\boldsymbol{\delta}})$.  It then follows that (1)~the efficient GMM
estimator is OLS (which actually is true without conditional homoskedasticity as long as
$\mathbf{z}_i = \mathbf{x}_i$), (2)~the restricted efficient GMM estimator subject to the constraints
of the null hypothesis is the restricted OLS (whose objective function is not $J$ but
$SSR$), and (3)~the Wald statistic, which is numerically equal to the $LR$ statistic, can
be calculated as the difference in $SSR$ with and without the imposition of the null,
normalized to $\hat{\sigma}^2$.  This last result confirms the derivation in Section~2.6
of the Wald statistic by the Likelihood-Ratio principle.

%% ─── Testing a Subset of Orthogonality Conditions ──────────────────────────
\subsection*{Testing a Subset of Orthogonality Conditions}

In Section~3.6 we introduced the statistic $C$ for testing a subset of orthogonality
conditions.  It utilizes two efficient GMM estimators of the same equation, one
using the full set $\mathbf{x}_i$ of instruments and the other using only a subset, $\mathbf{x}_{i1}$,
of $\mathbf{x}_i$.  To examine what expression it takes under conditional homoskedasticity, let
$\mathbf{X}_1$ $(n \times K_1)$ be the data matrix whose $i$-th row is $\mathbf{x}_{i1}'$.  Because the
two-step efficient GMM estimator is the 2SLS estimator under conditional homoskedasticity,
the two GMM estimators are given by
\begin{align}
  \hat{\boldsymbol{\delta}}
  &= (\mathbf{Z}'\mathbf{P}\mathbf{Z})^{-1}\, \mathbf{Z}'\mathbf{P}\mathbf{y}
  \quad \text{with } \mathbf{P}
  = \mathbf{X}(\mathbf{X}'\mathbf{X})^{-1}\mathbf{X}',
  \label{eq:3.8.15} \\
  \tilde{\boldsymbol{\delta}}
  &= (\mathbf{Z}'\mathbf{P}_1\mathbf{Z})^{-1}\, \mathbf{Z}'\mathbf{P}_1\mathbf{y}
  \quad \text{with } \mathbf{P}_1
  = \mathbf{X}_1(\mathbf{X}_1'\mathbf{X}_1)^{-1}\mathbf{X}_1'.
  \label{eq:3.8.16}
\end{align}
And the $C$ statistic becomes the difference in two Sargan statistics:
\begin{equation}
  C = \frac{\hat{\boldsymbol{\varepsilon}}'\mathbf{P}\hat{\boldsymbol{\varepsilon}}
    - \tilde{\boldsymbol{\varepsilon}}'\mathbf{P}_1\tilde{\boldsymbol{\varepsilon}}}
    {\hat{\sigma}^2},
  \label{eq:3.8.17}
\end{equation}
where
\[
  \hat{\boldsymbol{\varepsilon}} \equiv \mathbf{y}
  - \mathbf{Z}\hat{\boldsymbol{\delta}}, \quad
  \tilde{\boldsymbol{\varepsilon}} \equiv \mathbf{y}
  - \mathbf{Z}\tilde{\boldsymbol{\delta}}, \quad
  \hat{\sigma}^2 \equiv \frac{\hat{\boldsymbol{\varepsilon}}'\hat{\boldsymbol{\varepsilon}}}{n}.
\]
As seen for the case without conditional homoskedasticity, $C$ is guaranteed to be
nonnegative in finite samples if the same matrix $\hat{\mathbf{S}}$ is used throughout, which under
conditional homoskedasticity amounts to using the same estimate of the error variance,
$\hat{\sigma}^2$, to deflate both
$\hat{\boldsymbol{\varepsilon}}'\mathbf{P}\hat{\boldsymbol{\varepsilon}}$ and
$\tilde{\boldsymbol{\varepsilon}}'\mathbf{P}_1\tilde{\boldsymbol{\varepsilon}}$, as in
\eqref{eq:3.8.17}.  By Proposition~3.7, this $C$ is asymptotically $\chi^2(K - K_1)$.

Perhaps more popular is the test statistic proposed by Hausman and Taylor
(1980) and further examined by Newey (1985).  $\hat{\boldsymbol{\delta}}$ is asymptotically more efficient
than $\tilde{\boldsymbol{\delta}}$ because $\hat{\boldsymbol{\delta}}$ exploits more orthogonality conditions.
Therefore, $\avar(\tilde{\boldsymbol{\delta}}) \geq \avar(\hat{\boldsymbol{\delta}})$.
Furthermore, as you will be asked to show (Analytical Exercise~9), under the same
hypothesis guaranteeing $C$ to be asymptotically chi-squared,
\begin{equation}
  \avar(\tilde{\boldsymbol{\delta}} - \hat{\boldsymbol{\delta}})
  = \avar(\tilde{\boldsymbol{\delta}}) - \avar(\hat{\boldsymbol{\delta}}).
  \label{eq:3.8.18}
\end{equation}
(If you have been exposed to the Hausman test in the ML context, you can recognize
this as the GMM version of the Hausman principle.)  By Proposition~3.9,
$\avar(\hat{\boldsymbol{\delta}})$ is consistently estimated by
\begin{equation}
  \widehat{\avar(\hat{\boldsymbol{\delta}})}
  \equiv n \cdot \hat{\sigma}^2 \cdot (\mathbf{Z}'\mathbf{P}\mathbf{Z})^{-1}.
  \label{eq:3.8.19}
\end{equation}
Similarly, a consistent estimator of $\avar(\tilde{\boldsymbol{\delta}})$ is
\begin{equation}
  \widehat{\avar(\tilde{\boldsymbol{\delta}})}
  \equiv n \cdot \hat{\sigma}^2 \cdot (\mathbf{Z}'\mathbf{P}_1\mathbf{Z})^{-1}.
  \label{eq:3.8.20}
\end{equation}
Here, as in the calculation of the $C$ statistic, the same estimate, $\hat{\sigma}^2$, is used
throughout, in order to guarantee the test statistic below to be nonnegative.  The
resulting estimator of $\avar(\tilde{\boldsymbol{\delta}} - \hat{\boldsymbol{\delta}})$ is
\begin{equation}
  \widehat{\avar(\tilde{\boldsymbol{\delta}} - \hat{\boldsymbol{\delta}})}
  = \hat{\sigma}^2 \cdot n \cdot
    \bigl[(\mathbf{Z}'\mathbf{P}_1\mathbf{Z})^{-1}
    - (\mathbf{Z}'\mathbf{P}\mathbf{Z})^{-1}\bigr].
  \label{eq:3.8.21}
\end{equation}
Hausman and Taylor (1980) have shown that (1)~this matrix in finite samples is
positive semidefinite (nonnegative definite) but not necessarily nonsingular, but
(2)~for any generalized inverse\footnote{A \emph{generalized inverse}, $\mathbf{A}^-$, of a
matrix $\mathbf{A}$ is any matrix satisfying $\mathbf{A}\mathbf{A}^-\mathbf{A} = \mathbf{A}$.  If
$\mathbf{A}$ is square and nonsingular, then $\mathbf{A}^-$ is unique and equal to $\mathbf{A}^{-1}$.}
of this matrix, the \textbf{Hausman statistic}
\begin{align}
  H &\equiv \sqrt{n}\,(\tilde{\boldsymbol{\delta}} - \hat{\boldsymbol{\delta}})'\,
  \bigl\{\hat{\sigma}^2 \cdot n \cdot
  [(\mathbf{Z}'\mathbf{P}_1\mathbf{Z})^{-1}
  - (\mathbf{Z}'\mathbf{P}\mathbf{Z})^{-1}]\bigr\}^{-}\,
  \sqrt{n}\,(\tilde{\boldsymbol{\delta}} - \hat{\boldsymbol{\delta}}) \notag \\
  &= \frac{(\tilde{\boldsymbol{\delta}} - \hat{\boldsymbol{\delta}})'\,
  [(\mathbf{Z}'\mathbf{P}_1\mathbf{Z})^{-1}
  - (\mathbf{Z}'\mathbf{P}\mathbf{Z})^{-1}]^{-}\,
  (\tilde{\boldsymbol{\delta}} - \hat{\boldsymbol{\delta}})}{\hat{\sigma}^2}
  \label{eq:3.8.22}
\end{align}
is invariant to the choice of the generalized inverse and is asymptotically chi-squared
with $\min(K - K_1, L - s)$ degrees of freedom, where
\[
  s = \#\mathbf{z}_i \cap \mathbf{x}_{i1}
  = \text{number of regressors which are retained as instruments in } \mathbf{x}_{i1}.
\]

What is the relationship between $C$ and $H$ under conditional homoskedasticity?
It can be shown (see Newey, 1985) that:

If $K - K_1 \leq L - s$ so that both $H$ and $C$ have the same degrees of freedom, then
$H = C$ (numerically equal).  Otherwise, the two statistics are numerically different
and have different degrees of freedom.

One frequent case where $K - K_1 \leq L - s$ holds is when $\mathbf{x}_{i2}$ is a subset of
regressors, that is, when the suspect instruments are a subset of regressors.  In this case,
the formulas \eqref{eq:3.8.17} and \eqref{eq:3.8.22} are two alternative ways of calculating
numerically the same statistic.  In the case where $K - K_1 > L - s$, the degrees of freedom
for $H$ are less than those for the $C$ statistic $(K - K_1)$.  For this reason, the Hausman
test, unlike the $C$ test, is not consistent against some local alternatives.\footnote{The
Hausman statistic can be generalized to the case of conditional heteroskedasticity, but it
is not practical because the degrees of freedom depend on the unknown values of the
matrices $\boldsymbol{\Sigma}_{\mathbf{xz}}$ and $\mathbf{S}$.  See Newey (1985).}

%% ─── Testing Conditional Homoskedasticity ──────────────────────────────────
\subsection*{Testing Conditional Homoskedasticity}

For the OLS case with predetermined regressors, as shown in Section~2.7, there
is a convenient $nR^2$ test of conditional homoskedasticity.  In the present case, the
$nR^2$ statistic obtained by regressing the squared 2SLS residuals on a constant and
second-order cross products of the instrumental variables turns out \emph{not} to have the
desired asymptotic distribution.  A test statistic that is asymptotically chi-squared
is available but is extremely cumbersome.  See White (1982, note~2).

%% ─── Testing for Serial Correlation ────────────────────────────────────────
\subsection*{Testing for Serial Correlation}

For the OLS case, we developed in Section~2.10 tests for serial correlation in the
error term.  More specifically, under (2.10.15) (which is stronger than Assumption~2.3
or~3.3) and (2.10.16) (which is stronger than Assumption~2.7 or~3.7 of conditional
homoskedasticity), the modified Box--Pierce $Q$ given in (2.10.20) can be used for
testing the null of no serial correlation in the error term, and this statistic is
asymptotically equivalent (in the stronger sense of the plim of the difference being
zero) to the $nR^2$ statistic from regression (2.10.21).  Can this test be extended to the
case where the regressors $\mathbf{z}_i$ are endogenous?  If the instruments $\mathbf{x}_i$ satisfy
(2.10.15) and (2.10.16), then the argument in Appendix~2.B of Chapter~2 can be
generalized to produce a modified $Q$ statistic that is asymptotically chi-squared under
the null of no serial correlation.  However, the expression for the statistic is more
complicated than (2.10.20) and so is not presented here.  This modified $Q$ statistic for
the case of endogenous regressors does not seem to be asymptotically equivalent to
the $nR^2$ statistic from regression (2.10.21).

%% ─── Questions for Review ──────────────────────────────────────────────────
\bigskip
\begin{center}\rule{0.6\textwidth}{0.4pt}\end{center}
\noindent\textsc{Questions for Review}
\medskip

\begin{enumerate}
\item (GMM with conditional homoskedasticity) \; In efficient two-step GMM estimation,
  $\hat{\mathbf{S}}$ in the second step is calculated from the first step by \eqref{eq:3.5.10} using
  the first-step residuals.  Under conditional homoskedasticity, derive the asymptotic
  variance of the two-step estimator.  Is it the same as \eqref{eq:3.8.4}?
  \textbf{Hint:} Under conditional homoskedasticity,
  \eqref{eq:3.5.10} $\pto \sigma^2 \boldsymbol{\Sigma}_{\mathbf{xx}}$.

\item (2SLS without conditional homoskedasticity) \; Is the 2SLS consistent when
  conditional homoskedasticity does not hold?  Derive the plim of \eqref{eq:3.8.2} and
  $\avar(\hat{\boldsymbol{\delta}}_{\text{2SLS}})$ without assuming conditional homoskedasticity.
  Is the 2SLS as efficient as the two-step efficient GMM (i.e., is its asymptotic variance
  as small as \eqref{eq:3.5.13})?  \textbf{Hint:} 2SLS is a GMM estimator with a choice of
  $\hat{\mathbf{W}}$ that is not necessarily efficient without conditional homoskedasticity.

\item (GLS interpretation of 2SLS) \; Verify that the 2SLS estimator can be written as
  a GLS estimator if $\mathbf{S}_{\mathbf{xz}}$ and $\mathbf{s}_{\mathbf{xy}}$ are interpreted as
  the data matrix of regressors and the data vector of the dependent variable and
  $\mathbf{S}_{\mathbf{xx}}$ as the variance matrix of the error term.

\item Provide two cases in which $\hat{\boldsymbol{\delta}}_{\text{2SLS}}$ and
  $\hat{\boldsymbol{\delta}}(\hat{\mathbf{S}}^{-1})$ are asymptotically equivalent in the sense that
  \[
    \sqrt{n}\,[\hat{\boldsymbol{\delta}}_{\text{2SLS}}
    - \hat{\boldsymbol{\delta}}(\hat{\mathbf{S}}^{-1})] \pto \mathbf{0}.
  \]
  \textbf{Hint:} Keywords are ``conditional homoskedasticity'' and ``exactly identified.''

\item Suppose the equation is just identified.  Show that 2SLS \eqref{eq:3.8.3p} reduces to IV
  \eqref{eq:3.4.4}.

\item (Sargan as $nR^2$) \; Prove that Sargan's statistic \eqref{eq:3.8.10p} equals $nR_{uc}^2$,
  where $R_{uc}^2$ is the uncentered $R$-squared from a regression of
  $\hat{\boldsymbol{\varepsilon}}$ on $\mathbf{X}$.
  \textbf{Hint:} Review Question~8 of Section~1.2.

\item (When $\mathbf{z}_i$ is a strict subset of $\mathbf{x}_i$) \; Suppose $\mathbf{z}_i$ is
  a strict subset of $\mathbf{x}_i$.  So $\mathbf{x}_i$ includes, in addition to the regressors (which
  are all predetermined), some other predetermined variables.  We have shown that the
  efficient GMM estimator is OLS under conditional homoskedasticity.  Does the result
  remain true without conditional homoskedasticity?  \textbf{Hint:} Review Question~5 of
  Section~3.5.
\end{enumerate}
